% =========================================================
%  config/styles.tex
%  Warna, tipe kolom, format judul, header/footer, keterangan tabel
%  Skema warna dan tipe kolom diambil dari rps/config/styles.tex
%  agar tampilan buku ajar dan RPS satu keluarga.
% =========================================================

% --- Warna institusi ------------------------------------
\definecolor{ubPrimary}{HTML}{1F3864}
\definecolor{ubAccent}{HTML}{2E75B6}
\definecolor{ubHead}{HTML}{D9E2F3}
\definecolor{ubAlt}{HTML}{F2F6FC}
\definecolor{ubLine}{HTML}{BFBFBF}

% --- Warna semantik untuk kotak & kode ------------------
\definecolor{ubGreen}{HTML}{1E7B45}
\definecolor{ubGreenBg}{HTML}{F1F8F2}
\definecolor{ubRed}{HTML}{B3261E}
\definecolor{ubRedBg}{HTML}{FDF2F1}
\definecolor{ubWarn}{HTML}{A35A05}
\definecolor{ubWarnBg}{HTML}{FFF8EC}
\definecolor{ubGray}{HTML}{F6F6F6}

% --- Tipe kolom -----------------------------------------
\newcolumntype{L}[1]{>{\RaggedRight\arraybackslash}p{#1}}
\newcolumntype{C}[1]{>{\Centering\arraybackslash}p{#1}}
\newcolumntype{J}[1]{>{\justifying\arraybackslash}p{#1}}
\newcolumntype{Y}{>{\RaggedRight\arraybackslash}X}

% --- Jarak antar paragraf -------------------------------
\setlength{\parindent}{0pt}
\setlength{\parskip}{0.5em}
\renewcommand{\arraystretch}{1.25}

% --- Toleransi pemenggalan baris ------------------------
\setlength{\emergencystretch}{2em}
\tolerance=1500
\hyphenpenalty=200
\widowpenalty=1000
\clubpenalty=1000

% --- Nama bagian buku (bahasa Indonesia) ----------------
% babel menetapkan nama bagian pada saat \begin{document}, sehingga
% \renewcommand biasa pada preamble akan ditimpa. Karena itu nama bagian
% diubah melalui mekanisme babel (\addto\captionsindonesian).
\addto\captionsindonesian{%
  \renewcommand{\bibname}{Daftar Pustaka}%
  \renewcommand{\chaptername}{Bab}%
  \renewcommand{\contentsname}{Daftar Isi}%
  \renewcommand{\listfigurename}{Daftar Gambar}%
  \renewcommand{\listtablename}{Daftar Tabel}%
  \renewcommand{\figurename}{Gambar}%
  \renewcommand{\tablename}{Tabel}%
  \renewcommand{\appendixname}{Lampiran}%
}

% --- Format judul bab -----------------------------------
\titleformat{\chapter}[display]
  {\normalfont\LARGE\bfseries\color{ubPrimary}}
  {\filright\normalsize\bfseries\color{ubAccent}\MakeUppercase{\chaptertitlename}~\thechapter}
  {0.7ex}
  {\Large}
  [\vspace{0.5ex}{\color{ubAccent}\titlerule[0.9pt]}]
\titlespacing*{\chapter}{0pt}{6pt}{26pt}

% --- Format judul section -------------------------------
\titleformat{\section}
  {\normalfont\large\bfseries\color{ubPrimary}}
  {\thesection}{0.5em}{}
\titlespacing*{\section}{0pt}{1.5em}{0.6em}

\titleformat{\subsection}
  {\normalfont\normalsize\bfseries\color{ubAccent}}
  {\thesubsection}{0.5em}{}
\titlespacing*{\subsection}{0pt}{1.1em}{0.4em}

\titleformat{\subsubsection}
  {\normalfont\normalsize\bfseries}
  {\thesubsubsection}{0.5em}{}
\titlespacing*{\subsubsection}{0pt}{0.9em}{0.3em}

% --- Format judul bab tak bernomor ----------------------
\titleformat{name=\chapter,numberless}[display]
  {\normalfont\LARGE\bfseries\color{ubPrimary}}
  {}
  {0pt}
  {\Large}
  [\vspace{0.5ex}{\color{ubAccent}\titlerule[0.9pt]}]
\titlespacing*{name=\chapter,numberless}{0pt}{6pt}{26pt}

% --- Header & footer ------------------------------------
% Front matter memakai gaya polos (nomor halaman di tengah bawah);
% gaya fancy dengan kepala halaman dipakai mulai main matter.
\pagestyle{plain}
\fancypagestyle{mainmatter}{%
  \fancyhf{}%
  \fancyhead[L]{\footnotesize\color{ubPrimary}\NamaPT}%
  \fancyhead[R]{\footnotesize\color{ubPrimary}\leftmark}%
  \fancyfoot[C]{\footnotesize\thepage}%
  \renewcommand{\headrulewidth}{0.4pt}%
  \renewcommand{\footrulewidth}{0pt}%
}
% \@chapapp bernilai "Bab" pada bagian utama dan berubah menjadi "Lampiran"
% setelah \appendix. Memakai \chaptername secara langsung akan membuat judul
% berjalan pada lampiran tertulis "Bab B", sebab \chaptername tidak berubah.
\makeatletter
\renewcommand{\chaptermark}[1]{\markboth{\@chapapp\ \thechapter.\ #1}{}}
\makeatother
\renewcommand{\sectionmark}[1]{\markright{\thesection\ #1}}

% --- Keterangan tabel & gambar --------------------------
\captionsetup{
  font=small,
  labelfont={bf,color=ubPrimary},
  skip=4pt
}

% --- Daftar (list) --------------------------------------
\setlist[itemize]{topsep=2pt,itemsep=1pt,parsep=0pt,leftmargin=1.3em}
\setlist[enumerate]{topsep=2pt,itemsep=2pt,parsep=0pt,leftmargin=1.5em}

% Daftar khusus untuk checklist kompetensi
\newlist{checklist}{itemize}{1}
\setlist[checklist]{label=$\square$,leftmargin=1.7em,itemsep=3pt,topsep=3pt}

% Daftar khusus untuk langkah praktikum
\newlist{langkah}{enumerate}{1}
\setlist[langkah]{label=\textbf{\arabic*.},leftmargin=1.9em,itemsep=3pt,topsep=3pt}

% --- Tabel ringkas (keterangan di bawah tabel) ----------
\newcommand{\tabsource}[1]{%
  \par\vspace{2pt}{\footnotesize\itshape #1\par}%
}

% --- Garis tabel (booktabs) -----------------------------
% Header tabel ditandai garis tebal di atas dan tipis di bawah,
% sebagai pengganti pewarnaan latar baris (colortbl tidak tersedia).
\renewcommand{\toprule}{\specialrule{1.3pt}{0pt}{0pt}}
\renewcommand{\midrule}{\specialrule{0.5pt}{0pt}{0pt}}

% --- Daftar isi (tocloft) -------------------------------
\renewcommand{\cftchapfont}{\bfseries\color{ubPrimary}}
\renewcommand{\cftchappagefont}{\bfseries\color{ubPrimary}}
\renewcommand{\cftsecfont}{\color{black}}
\renewcommand{\cfttoctitlefont}{\hfill\Large\bfseries\color{ubPrimary}}
\renewcommand{\cftaftertoctitle}{\hfill}
\renewcommand{\cftloftitlefont}{\hfill\Large\bfseries\color{ubPrimary}}
\renewcommand{\cftafterloftitle}{\hfill}
\renewcommand{\cftlottitlefont}{\hfill\Large\bfseries\color{ubPrimary}}
\renewcommand{\cftafterlottitle}{\hfill}

% --- Tampilan tautan ------------------------------------
\hypersetup{
  colorlinks = true,
  linkcolor  = ubPrimary,
  urlcolor   = ubAccent,
  citecolor  = ubAccent,
  filecolor  = ubAccent
}
